\begin{afterword}
\summaryhead

The post contrasts two images. Lady Justice carries scales, on which whatever pulls one side down pushes the other up. People treat debate the same way, like a sport in which every point against one side counts for the other.

A study of the affect heuristic asked whether people merge a technology's risks and benefits into one good or bad feeling. The post's example is a reactor that makes less waste but is more likely to melt down. The two facts may be cited by opposite sides, but they are separate physical questions with separate answers. People who judge by overall feeling, the post says, will rate the meltdown risk lower after hearing about the low waste, and so get a factual question wrong.

Scales are appropriate, following E. T. Jaynes, only for a single factual question with two possible answers, such as whether someone is guilty. For anything more complex, Lady Rationality keeps a notebook of facts that are on no one's side.

\responsehead

The post's central distinction is correct and useful. Two facts can be cited by opposite sides of a debate and still be independent, so learning one should not change your estimate of the other. The study it cites supports the claim that people do this, and the study's authors give the reason it counts as an error: in the world, risks and benefits tend to go together, while in people's minds they run opposite. The appeal to Jaynes also checks out. With two hypotheses, evidence moves the odds between them, and a single balance is the right picture.

Two things are weaker than they look. The first is the example. The post says that if you tell people a reactor produces less waste, ``they rate its probability of meltdown as lower.'' The study did not test waste against meltdown. It gave people a paragraph about the risks or the benefits of nuclear power in general and measured how their ratings of both changed. The result runs in the same direction, but the post presents its own illustration as a finding.

The second is the closing image. The post ends by saying that for any issue more complex than a single factual question, Justice ``should relinquish her scales or relinquish her sword.'' But the notebook only settles the facts. A decision about a reactor, or about the banned-products shops of ``Policy Debates Should Not Appear One-Sided,'' still has to weigh the facts against each other, and that post's example of real tough-mindedness was keeping the shops open ``because we did this cost-benefit calculation.'' The post's argument supports keeping the facts apart while estimating them. It does not show that weighing them is a distortion.

In short: a correct and useful distinction between facts and sides, supported by its study, illustrated with an example the study did not test, and closed with an image that claims more than the argument shows.
\end{afterword}
